\subsection{Geometria}

\begin{small}
\begin{itemize}
    \item \textbf{Fórmula de Euler:} Em um grafo planar ou poliedro convexo, temos:
    $ V - E + F = 2 $
    onde $V$ é o número de vértices, $E$ o número de arestas e $F$ o número de faces.

    \item \textbf{Teorema de Pick:} Para polígonos com vértices em coordenadas inteiras:
    \[
    \text{Área} = i + \frac{b}{2} - 1
    \]
    onde $i$ é o número de pontos interiores e $b$ o número de pontos sobre o perímetro.

    \item \textbf{Teorema das Duas Orelhas (Two Ears Theorem):} Todo polígono simples com mais de três vértices possui pelo menos duas "orelhas" — vértices que podem ser removidos sem gerar interseções. A remoção repetida das orelhas resulta em uma triangulação do polígono.

    \item \textbf{Incentro de um Triângulo:} É o ponto de interseção das bissetrizes internas e centro da circunferência inscrita. Se $a$, $b$ e $c$ são os comprimentos dos lados opostos aos vértices $A(X_a, Y_a)$, $B(X_b, Y_b)$ e $C(X_c, Y_c)$, então o incentro $(X, Y)$ é dado por:
    \[
    X = \frac{aX_a + bX_b + cX_c}{a + b + c}, \quad
    Y = \frac{aY_a + bY_b + cY_c}{a + b + c}
    \]

    \item \textbf{Triangulação de Delaunay:} Uma triangulação de um conjunto de pontos no plano tal que nenhum ponto está dentro do círculo circunscrito de qualquer triângulo. Essa triangulação:
    \begin{itemize}
        \item Maximiza o menor ângulo entre todos os triângulos.
        \item Contém a árvore geradora mínima (MST) euclidiana como subconjunto.
    \end{itemize}

    \item \textbf{Fórmula de Brahmagupta:} Para calcular a área de um quadrilátero cíclico (todos os vértices sobre uma circunferência), com lados $a$, $b$, $c$ e $d$:
    \[
    s = \frac{a + b + c + d}{2}, \quad
    \text{Área} = \sqrt{(s - a)(s - b)(s - c)(s - d)}
    \]
    Se $d = 0$ (ou seja, um triângulo), ela se reduz à fórmula de Heron:
    \[
    \text{Área} = \sqrt{(s - a)(s - b)(s - c)s}
    \]
\end{itemize}
\end{small}

\subsection{Point in Polygon e Truques Geométricos}

\begin{small}
\begin{itemize}
    \item \textbf{Ray Casting (Curva de Jordan):} Resolve ponto em polígono geral em $O(N)$. Traça-se um raio do ponto $P$ até o infinito. Se cruzar as arestas um número ímpar de vezes, $P$ está dentro.
    \textbf{Truque dos Coprimos:} Para evitar \textit{edge cases} (raio bater num vértice ou aresta colinear), use um vetor de direção com primos altos:
    \[
    P_{\infty} = (P.x + 10^9+7, P.y + 10^9+9)
    \]
    Isso inclina a reta o suficiente para garantir que ela nunca acerte uma coordenada inteira exata, reduzindo o problema a uma simples checagem de interseção de segmentos.

    \item \textbf{Winding Number (Soma dos Ângulos):} Alternativa em $O(N)$ baseada na soma dos ângulos entre $P$ e vértices adjacentes:
    \[
    \sum \text{Ângulos} = \begin{cases} \pm 2\pi, & \text{dentro} \\ 0, & \text{fora} \end{cases}
    \]
    \textbf{Dica:} Implemente as checagens usando apenas \textbf{Produto Vetorial} (verificando semiplanos) para manter a aritmética inteira e fugir da imprecisão do \texttt{atan2}.

    \item \textbf{Prefix Sum de Áreas (Shoelace):} Usado para responder \textit{queries} de área de sub-polígonos (cortes) em $O(1)$. Mantém-se o prefixo do produto vetorial:
    \[
    S[k] = S[k-1] + (V_{k-1} \times V_k)
    \]
    A área da fatia conectando $V_i$ a $V_j$ é:
    \[
    \text{Area}_{i,j} = \frac{1}{2} \Big| S[j] - S[i] + (V_j \times V_i) \Big|
    \]

    \item \textbf{Busca Binária em Polígono Convexo:} Resolve ponto em polígono estritamente convexo em $O(\log N)$.
    \begin{itemize}
        \item Fixe o vértice $V_0$.
        \item Use busca binária baseada no \texttt{ccw} (Cross Product) para achar o triângulo angular $(V_0, V_k, V_{k+1})$ que contém $P$.
        \item Verifique se $P$ está no semiplano interno da aresta $V_k \to V_{k+1}$.
    \end{itemize}
\end{itemize}
\end{small}

% credits: https://github.com/gabrielpessoa1/ICPC-Library/